\section{Methods and Benchmark Design}
\label{sec:methods}

% Question: What do the methods of this project consist of?
% Answer: Two parts: NOVA, the detector-level inference method, and the
% injector, which hides a known spectrum inside real JWST data so recovery
% accuracy can be measured absolutely.
% Evidence: The structure of this section.

The methods comprise two parts. NOVA is the inference method: a forward
model of the calibrated detector time series whose parameters are fitted
directly to the retained pixels. The injector provides the test: it hides
a known transmission spectrum inside real out-of-transit JWST data, so
that the accuracy of any recovery method can be measured absolutely. The
exact configuration evaluated in this assessment, and its provenance, are
specified in Appendix~\ref{app:nova}.

\subsection{\nova{}}

% Question: What problem does NOVA solve, and in what form?
% Answer: It inverts a forward model of the detector: the same physical
% quantities that generate the pixels are inferred from them in one
% likelihood.
% Evidence: Sections 1.3 and 1.4.

NOVA treats spectral recovery as a single inverse problem. A forward
model maps the physical description of the observation, the transmission
spectrum $D(\lambda)$, the limb darkening $u(\lambda)$, the stellar
continuum, the spatial response, and the additive background, to a
prediction for every retained detector sample; the inverse problem infers
the transmission spectrum, the limb darkening, the continuum, and the
background from the observed samples, conditional on the fixed response,
orbit, and source-curvature factor. Nothing is extracted, binned,
or fitted in a separate downstream stage: the wavelength-dependent
transit is estimated in the same likelihood that describes the detector.

\subsubsection{The detector forward model}

% Question: What does NOVA predict for each detector sample?
% Answer: Continuum times exposure-integrated transit times spatial
% response, corrected by a fixed source-curvature factor, plus an additive
% background.
% Evidence: Presentation slide 12 and the exact method note.

The retained detector samples are organised into wavelength groups: one
group $g$ collects the pixels $P_g$ of one spectral order that share one
native detector column, and $g=g(p)$ denotes the group of pixel $p$. For
integration $t$ and pixel $p$, NOVA predicts
\begin{equation}
 \widehat Y_{tp}
 = \gamma_t\, C_{tg}\,
   \mathcal{T}_{tg}(D_g, \bm{u}_g;\, \Omega)\,
   q^\star_{pg}
 + B_{tp},
 \label{eq:nova-forward}
\end{equation}
and, where order supports overlap, the source term is summed over the
groups $\mathcal{G}(p)$ contributing to the pixel. Reading outward from
the star: the continuum $C_{tg}$ sets the brightness of the group at time
$t$; the transit factor $\mathcal{T}_{tg}$ dims that light according to
the depth $D_g$ and limb darkening $\bm{u}_g$ at the group's wavelength;
the factor $\gamma_t$ applies a small fixed correction for slow source
variation over the visit; the spatial response $q^\star_{pg}$ distributes
the group's light across its pixels, with
$\sum_{p\in P_g} q^\star_{pg}=1$; and $B_{tp}$ is the additive
background. The group geometry, wavelength assignment, and response
support follow the current public trace and wavelength calibration
\citep{BainesEtAl2023Trace,BainesEtAl2023Wavelength} and the associated
spectral-response kernel
\citep{DarveauBernierEtAl2022,JWSTPipeline2025}. The background is
additive and is never multiplied by the transit or by $\gamma_t$, so a
change in transit depth cannot be imitated by dimming the background
along with the star.

Valid samples have finite values, positive uncertainties, and no
\texttt{DO\_NOT\_USE} flag; 97{,}557 pixels per integration are retained,
organised into 2{,}720 groups after six unreliable edge groups are
excluded.

\subsubsection{The shared spectrum and transit model}

% Question: How are the depths and the transit generated?
% Answer: One physical spectrum shared by both orders, an exact Keplerian
% transit averaged over each exposure, and small fitted limb-darkening
% deviations around a theory reference.
% Evidence: The exact method note and presentation slides 12 and 13.

Both spectral orders evaluate the same physical spectrum on their own
wavelength grids, $D_g=D(\lambda_g)$. The spectrum is carried by $K=160$
adaptive basis coefficients,
\begin{equation}
  D_j = D_{\mathrm{global}}\,
  \frac{\exp\!\left[(\bm{H}_v\bm{v})_j\right]}
       {\sum_k w_k \exp\!\left[(\bm{H}_v\bm{v})_k\right]} ,
  \label{eq:depth-basis}
\end{equation}
which separates the achromatic depth $D_{\mathrm{global}}$ from a
weighted-zero-mean chromatic morphology and keeps all depths positive.
The model also contains a two-coordinate order-discrepancy term $\Delta$,
with order 1 seeing $D(\lambda)+\Delta$ and order 2 seeing
$D(\lambda)-\Delta$; here it is pinned to zero by its numerical bounds
(Appendix~\ref{app:nova}), so the two orders share one spectrum exactly.
No atmospheric template, wavelength-local smoother, or spectral prior is
imposed.

The orbital state
$\Omega=(P,\,a/R_\star,\,e,\,\omega_{\mathrm{peri}},\,b,\,t_0)$ is held
fixed at the values listed in Appendix~\ref{app:nova}. The instantaneous
limb-darkened transit is computed with an exact Keplerian engine
\citep{Jaxoplanet2025}, and $\mathcal{T}_{tg}$ is its average over the
finite exposure, evaluated with 21 weighted samples per integration.
Quadratic limb darkening \citep{MandelAgol2002} uses the bounded
coordinates of \citet{Kipping2013},
\begin{equation}
 u_1=2\sqrt{q_1}\,q_2, \qquad
 u_2=\sqrt{q_1}\,(1-2q_2),
 \label{eq:kipping-coordinates}
\end{equation}
with the wavelength-dependent theory reference taken from the STAGGER
grid \citep{MagicEtAl2015}. The fitted freedom around that reference is
small: one offset and one normalised log-wavelength slope per transformed
coordinate and order, eight nonlinear coordinates in total, constrained
by theory-deviation and curvature penalties.

\subsubsection{The empirical spatial response}

% Question: Where does q-star come from?
% Answer: The transit itself reveals where the starlight lands: a robust
% per-pixel regression on a depth-free transit template, with amplitudes
% normalised within each group.
% Evidence: Presentation slide 14 and the exact method note.

The spatial response is estimated from the data themselves, using the
transit as its own probe: starlight dims during transit and the additive
background does not, so the transit-correlated amplitude of each pixel
measures how much starlight lands there. For each order, a depth-free
event template $\phi_o(t)$ is built from a robust median of
baseline-normalised bright-pixel residuals, and every qualifying pixel is
fitted robustly as
\begin{equation}
 Y_{tp} = A_{pg}\,\phi_{o}(t) + m_p + n_p\,\tau_t + \varepsilon_{tp},
 \label{eq:qstar-fit}
\end{equation}
with a level $m_p$, a linear drift $n_p\tau_t$, and the transit-correlated
amplitude $A_{pg}$. The response is the normalised amplitude,
$q^\star_{pg}=A_{pg}/\sum_{p'\in P_g}A_{p'g}$, so total group brightness
remains with the continuum. Two alternating folds, stratified across the
transit phases, provide independent estimates, and the response is
accepted only if the two folds agree in flux distribution and centroid
(Appendix~\ref{app:nova}). The adopted response has a known limitation:
its smoothing operates on an integer trace-relative coordinate, which
produces occasional centroid resets and a small amplitude excess at the
reddest columns. Because the continuum sets each group's brightness while
$q^\star$ sets its spatial profile, such an excess acts like a dilution
and biases the recovered transit shallow at the affected columns; its
impact is quantified in the Results.

\subsubsection{The continuum and the additive background}

% Question: How are the stellar continuum and the background modelled and
% calibrated?
% Answer: A level and slope per group for the continuum; eight spatial maps
% crossed with two temporal patterns for the background, selected and
% anchored on off-trace pixels.
% Evidence: Presentation slide 15 and the exact method note.

The continuum of each group is a linear function of time,
$C_{tg}=\beta_{0g}+\beta_{1g}\tau_t$, giving 5{,}440 linear coefficients.
The background is a sum of eight fixed spatial maps $\Phi_{pk}$ with two
temporal patterns,
\begin{equation}
 B_{tp}=\sum_{k=1}^{8}\Phi_{pk}
 \left[\alpha_{k0}+\alpha_{k1}\,G_1(t)\right],
 \label{eq:background-model}
\end{equation}
sixteen coefficients $\alpha_{km}$ in all. This model is calibrated
entirely on pixels away from the stellar traces. The spatial family and
rank were chosen by hold-out prediction on off-trace pixels, and $G_1(t)$
is the leading principal component of their temporal behaviour whose
absolute correlation with a transit-shaped reference stays below 0.25,
which limits the background's ability to absorb transit-shaped
variation. The off-trace fit also yields
an expectation $\bm{\alpha}_{\mathrm{off}}$ with precision
$\bm{N}_{\mathrm{off}}$, which enters the objective as a soft anchor: the
background is anchored, not fixed, and is refitted jointly with the
continuum at every step. No fixed background is subtracted at any stage.
The likelihood contains no Gaussian-process term for time-correlated
residuals; outliers are handled by the robust weights below, and residual
temporal correlation is tracked as a diagnostic.

\subsubsection{The source-curvature factor and detector uncertainties}

% Question: What are gamma_t and sigma_tp?
% Answer: A fixed out-of-transit curvature correction selected on held-out
% folds, and per-sample uncertainties scaled per order and per detector row
% from out-of-transit cross-validation.
% Evidence: The exact method note, Sections 8 and 10.

The factor $\gamma_t$ corrects for slow variation of the source over the
visit. It was calibrated on the out-of-transit integrations, split into
two folds: stable trace flux was integrated per order, the ratio of a
weighted quadratic in time to its affine approximation defined each
candidate, and null, common, and per-order candidates were compared by
held-out chi-squared. The common candidate was selected, because the
per-order one did not improve both orders individually
(Appendix~\ref{app:nova}). $\gamma_t$ has unit out-of-transit mean,
multiplies only the stellar term of Equation~\ref{eq:nova-forward}, and
adds no fitted parameters.

The per-sample uncertainties are
$\sigma_{tp}=\mathrm{ERR}_{tp}\,N_o\,s_{r(p)}$: the calibrated ERR array,
scaled once per order by $N_o$ and once per detector row by a factor
$s_{r(p)}\geq 1$ from bidirectional out-of-transit cross-validation, so
rows whose quiet-time residuals exceed their nominal uncertainties are
down-weighted. Both scales derive from out-of-transit data alone.

\subsubsection{The complete objective}

% Question: What does NOVA minimise?
% Answer: A robust detector mismatch over the retained mask plus the
% off-trace background anchor and the limb-darkening and order-discrepancy
% penalties.
% Evidence: Presentation slide 16.

With the standardised residual
\begin{equation}
 r_{tp}=\frac{Y_{tp}-\widehat Y_{tp}
 (D,\bm{u},\bm{\beta},\bm{\alpha},\Delta;\,q^\star,\Omega)}{\sigma_{tp}},
 \label{eq:residual}
\end{equation}
in which fitted parameters appear before the semicolon and fixed inputs
after it, the estimate is
\begin{equation}
\begin{split}
 &(\widehat D,\widehat{\bm{u}},\widehat{\bm{\beta}},
  \widehat{\bm{\alpha}},\widehat\Delta)
 = \arg\min\Big\{ \\
 &\quad\sum_{(t,p)\in\mathcal{M}}\rho_{1.345}(r_{tp}) \\
 &\quad+\tfrac{1}{2}(\bm{\alpha}-\bm{\alpha}_{\mathrm{off}})^{\mathsf T}
   \bm{N}_{\mathrm{off}}(\bm{\alpha}-\bm{\alpha}_{\mathrm{off}}) \\
 &\quad+\tfrac{1}{2}\|\bm{r}_{\mathrm{LD,theory}}(\bm{u})\|_2^2
  +\tfrac{1}{2}\|\bm{r}_{\mathrm{LD,curv}}(\bm{u})\|_2^2 \\
 &\quad+\tfrac{1}{2}\|\bm{r}_{\mathrm{order}}(\Delta)\|_2^2
 \Big\},
\end{split}
 \label{eq:objective}
\end{equation}
where $\mathcal{M}$ is the retained sample mask and $\rho_{1.345}$ is the
Huber loss: the detector mismatch, the off-trace background anchor, the
limb-darkening theory and smoothness penalties, and the order-discrepancy
penalty.

\subsubsection{The nested solver}

% Question: How is the objective minimised, and what counts as convergence?
% Answer: Robust weights outside, bounded trust-region steps in the middle,
% an exact linear solve inside; acceptance requires convergence gates and
% agreement of two independent starts.
% Evidence: Presentation slides 17 and 18, and the solver code.

The parameters split into a nonlinear block
$\bm{\theta}=(D,\bm{u},\Delta)$ and a linear block
$(\bm{\beta},\bm{\alpha})$, and the solver nests three layers. Innermost,
at fixed $\bm{\theta}$ and robust weights, variable projection
\citep{GolubPereyra1973} solves the linear problem exactly,
\begin{equation}
\begin{split}
 (\widehat{\bm{\beta}},\widehat{\bm{\alpha}})=
 \arg\min_{\bm{\beta},\bm{\alpha}}\;
 &\frac{1}{2}\sum_{(t,p)\in\mathcal{M}}\frac{w_{tp}}{\sigma_{tp}^2}
 \left[Y_{tp}-\widehat Y_{tp}\right]^2 \\
 &+\frac{1}{2}(\bm{\alpha}-\bm{\alpha}_{\mathrm{off}})^{\mathsf T}
   \bm{N}_{\mathrm{off}}(\bm{\alpha}-\bm{\alpha}_{\mathrm{off}}),
\end{split}
 \label{eq:varpro}
\end{equation}
by factorising the sparse continuum block and reducing the background
coupling to a sixteen-dimensional Schur complement; the profiled residual
and Jacobian differentiate through this exact inner optimum
(Appendix~\ref{app:nova}). In the middle, bounded trust-region reflective
least squares \citep{BranchColemanLi1999} minimises the profiled
objective over $\bm{\theta}$. Outermost, Huber reweighting
\citep{Huber1964} updates the weights,
$w_{tp}=\min(1,\,1.345/|e_{tp}|)$ for standardised residual $e_{tp}$,
holding them fixed during each trust-region solve.

A fit is accepted only when the weights and the objective have stopped
changing, the shared spectrum is stationary to within 0.01 and 0.05 ppm
on the two spectral summaries, the first-order optimality measures are
met, and a confirmation fit at the recomputed terminal weights
reproduces the solution; the remaining numerical gates are listed in
Appendix~\ref{app:nova}. Each analysis is started twice, independently:
once at the white-light depth with zero morphology and the
limb-darkening theory reference, and once with a fixed 250 ppm
morphology perturbation and small limb-darkening offsets. The two starts
must agree to within 0.1 and 0.5 ppm on the same summaries, and if both
converge, the one with the lower robust objective is selected.

\subsection{The injector}

% Question: How is a truth-known benchmark built from real JWST data?
% Answer: A known transit deficit, applied only to a denoised target-source
% expectation, is added to real out-of-transit integrations, so background
% and noise are real and the truth is known.
% Evidence: Presentation slide 20 and the injector record.

The injector answers the validation problem of
Section~\ref{sec:detector-to-validation}: real observations provide no
ground truth, so a truth is hidden inside real data. The baseline is the
observed out-of-transit detector time series of the WASP-17 NIRISS/SOSS
visit, 229 calibrated integrations that already contain the target star,
the additive background, field sources, and the realised noise. A chosen
transmission spectrum $D^{\mathrm{inj}}(\lambda)$ defines per-order
transit factors $\bm{T}^{\mathrm{true}}_{o,t}$, and the delivered cube is
\begin{equation}
\begin{split}
 Y^{\mathrm{inj}}_{tp}
 = {}& Y^{\mathrm{OOT}}_{tp} \\
 & + \sum_{o}\left[\bm{A}_o\!\left\{\bm{s}_{o,t}\odot
   \left(\bm{T}^{\mathrm{true}}_{o,t}-\bm{1}\right)\right\}\right]_p ,
\end{split}
 \label{eq:v45-injector}
\end{equation}
where $\bm{s}_{o,t}$ is a denoised expectation of the target source for
order $o$ and $\bm{A}_o$ is the corresponding detector projection, built
from the same public instrument references
\citep{BainesEtAl2023Trace,BainesEtAl2023Wavelength,DarveauBernierEtAl2022}.
Only the target-source expectation transits: the observed background,
contaminants, and noise are untouched, and the source expectation is
denoised so that no particular noise realisation is dimmed. Columns where
the target is bright but the source model lacks valid response support
raise an error rather than being silently clipped. The synthetic event
occupies integrations 51--177, leaving 102 event-complement integrations.
The fixed calibration masks are selected from this complement as detailed
in Appendix~\ref{app:nova}. The injected spectrum is revealed
only after a recovery's prediction has been fixed; until then, every
calibration and modelling choice rests on the data and the instrument
references alone (Appendix~\ref{app:benchmarking}).

\subsubsection{Comparison analyses and fairness}

% Question: Against what is NOVA compared, and on what terms?
% Answer: An authentic Ahsoka analysis of the same cube now, further
% pipelines later; fairness is defined at delivery, and the group-stage
% preprocessing experiment remains open.
% Evidence: The injector record and Louie et al. (2025).

Every recovery method receives the same cube, uncertainties, flags, time
grid, and wavelength support; unsupported bins are excluded
symmetrically, and each method keeps its own nuisance model and inference
strategy. The first external comparison is an Ahsoka reduction following
the analysis topology of the published WASP-17~b spectrum
\citep{LouieEtAl2025}. One fairness limitation is open: the principal
comparison pipelines apply their most distinctive $1/f$ correction on
detector groups before ramp fitting, and an integration-level cube
contains no groups. A paired experiment will therefore inject the same
deficit at group stage and at integration stage and run NOVA and the
comparison pipelines on both; it has not yet been carried out. Later
rounds will add Eureka!, supreme-SPOON/exoTEDRF, transitspectroscopy, and
the official JWST pipeline with ATOCA extraction
\citep{BellEtAl2022,Radica2024,Espinoza2022TransitSpectroscopy,%
JWSTPipeline2025,DarveauBernierEtAl2022,LouieEtAl2025}, each with its own
standard settings on the same delivered cube.

\subsubsection{Reporting support and evaluation metrics}

% Question: On what wavelength support are methods compared, and with which
% metrics?
% Answer: Whole bins of a common R=100 grid inside the directly calibrated
% wavelength domain, scored by absolute RMSE, mean bias, and morphology
% RMSE.
% Evidence: The exact method note and the benchmark appendix.

Recovered and injected spectra are compared on a common $R=100$ grid,
restricted to bins whose detector support lies entirely inside the
directly calibrated wavelength domain of the trace solution; beyond its
last training anchor the wavelength solution is a polynomial
extrapolation, and those bins are excluded in full. For this visit the
rule keeps 141 bins, the last ending near $2.740\,\mu\mathrm{m}$
(Appendix~\ref{app:benchmarking}); it is the same for every method and is
applied only after recovery.

With $e_i=\widehat D_i-D_i^{\mathrm{inj}}$ on the common grid, the scores
are
\begin{equation}
 \mathrm{RMSE}_{\mathrm{abs}}=
 \left(\frac{1}{N}\sum_i e_i^2\right)^{1/2},
 \label{eq:absolute-rmse}
\end{equation}
\begin{equation}
 \mathrm{RMSE}_{\mathrm{morph}}=
 \left[\frac{1}{N}\sum_i(e_i-\bar e)^2\right]^{1/2},
 \label{eq:morphology-rmse}
\end{equation}
so a constant offset is distinguished from an incorrect spectral shape,
together with mean and median bias, feature-amplitude error, and
order-specific residuals. Standardised residuals and the spectral
covariance from the profiled normal system accompany each recovery;
empirical coverage of injected depths requires repeated noise
realisations and is part of the planned ensemble. A low error is not
sufficient on its own: every score is read alongside the residual,
weight, conditioning, and convergence diagnostics of the fit that
produced it.
